\section{MetaShift-Bench audit framework}

\begin{figure}[tbp]
\centering
\includegraphics[width=\linewidth]{fig_audit_pipeline.pdf}
\caption{MetaShift-Bench evidence workflow. Public archives are canonicalized
before persistent metadata anchors and distinct physical donor candidates are
constructed. Stable regimes supply known synthetic truth, while every real
anchor receives a common comparison or a recorded failure reason. The diagram
is a workflow schematic; its arrows do not assert that metadata reveals a
physical instrument state.}
\label{fig:audit-pipeline}
\end{figure}

\subsection{Counterfactual construction}

For a complete anchor, the donor pool is ranked by pre-event correlation,
distance, paired days, and POC only after physical-site deduplication. The
nearest-neighbor difference-in-differences baseline places unit weight on the
top-ranked donor. Standard synthetic control solves the nonnegative simplex
problem
\[
  \min_{\substack{w_j\geq0\\\sum_jw_j=1}}
  \frac{1}{|T_{\mathrm{cal}}|}
  \sum_{t\in T_{\mathrm{cal}}}
  \left(z_{i,t}-\sum_{j\in N_a}w_jz_{j,t}\right)^2.
\]
The intercept is represented by calibration centering in the residual
definition of Section~2. All fitting data end before the 15-day buffer; no
post-anchor outcome is supplied to either method.

Fixed-prior MetaShift starts from the pre-event reliability preference
\[
  \pi_{ij}\propto
  \max(\rho_{ij},0)^2\exp(-d_{ij}/50),
\]
where $\rho_{ij}$ is the pre-event log-scale correlation and $d_{ij}$ is
distance in km. Its frozen real-audit objective is
\[
  \min_{\substack{w_j\geq0\\\sum_jw_j=1}}
  \operatorname{MSE}_{T_{\mathrm{cal}}}(z_i,Z_Nw)
  +0.1\sum_jw_j^2+0.1\sum_j(w_j-\pi_{ij})^2.
\]
This is a reproducible regularized weighting rule, not a learned instrument
taxonomy. In the synthetic benchmark only, cross-validated MetaShift tests a
frozen set of penalty pairs on the final 45 pre-event days, chooses the
smallest validation RMSE, and refits the selected pair on the supplied
pre-event calibration table. It never sees anchor-post data during selection.

\subsection{Single-series comparators}

The benchmark includes five transparent single-series baselines:
before-after median, Bayesian mean shift, CUSUM, rolling-MAD, and PELT.
Before-after median and Bayesian mean shift estimate a known-date 60-day
level change. CUSUM, rolling-MAD, and PELT score or locate changes in the
single target series \citep{killick2012,truong2020,fryzlewicz2014}. They
serve a necessary contrast: a method can detect an abrupt time-series pattern
without establishing that the pattern is local rather than regional.

\subsection{Interval and falsification diagnostics}

For fixed donor weights, the primary interval resamples pre- and post-anchor
residuals using circular moving blocks of seven days, with 1,000 fixed-seed
repetitions. It is conditional on previously chosen donors and weights; it
does not include selection uncertainty. The selection-aware nested procedure
re-block-resamples the calibration matrix, recomputes correlation eligibility,
selects three to five donors from the fixed candidate pool, refits the
nonnegative MetaShift weights, and separately resamples the pre/post residual
windows. It preserves temporal ordering and reports failed reselection or
effect replicates rather than silently replacing them.

The audit also computes stable post-transition time placebos,
donor-as-treated placebos, a date-resampling comparison, and
leave-one-donor-out refits. These procedures are falsification and robustness
diagnostics. They do not create random assignment, prove a physical mechanism,
or make a residual effect a measured instrument bias.

\subsection{Rule-based evidence synthesis and abstention}

\begin{table}[tbp]
\centering
\caption{Predeclared primary evidence-tier rules and abstention outcomes.}
\label{tab:evidence-tier-rules}
\scriptsize
\begin{tabular}{p{0.39\linewidth}rp{0.40\linewidth}}
\toprule
Required diagnostic or condition & Threshold & Audit outcome \\
\midrule
Complete common-method comparison & required & Missing input $\Rightarrow$ inconclusive \\
Selection-aware nested interval & required & Unavailable $\Rightarrow$ inconclusive \\
Unique stable post-transition time placebos & 50 & Fewer dates $\Rightarrow$ inconclusive \\
Raw time-placebo probability & 0.10 & Above cutoff $\Rightarrow$ not supported \\
BH-adjusted time-placebo q value & 0.10 & Unavailable $\Rightarrow$ inconclusive; above cutoff $\Rightarrow$ not supported \\
Fixed-weight conditional diagnostic interval & exclude 0 & Otherwise $\Rightarrow$ not supported \\
Leave-one-donor-out direction fraction & 0.90 & Below cutoff $\Rightarrow$ not supported \\
All available required conditions & pass & Supported candidate discontinuity \\
\bottomrule
\end{tabular}
\par\smallskip
\begin{minipage}{0.94\linewidth}
\footnotesize\textit{Note.} The rules synthesize observational diagnostics; they are not a supervised classifier, instrument-fault label, or physical-causality test. The nested interval is required for tier assignment but is not claimed to have synthetic coverage calibration.
\end{minipage}
\end{table}


The tier logic distinguishes unavailable evidence from failed available
diagnostics. An event is inconclusive if it lacks a common comparison,
selection-aware interval, sufficient time placebos, an adjusted placebo
quantity, or donor sensitivity. It is not supported when those required
inputs are available but a pre-event fit, fixed diagnostic interval, placebo
screen, or donor-direction condition fails. A supported candidate satisfies
all available required conditions. Benjamini--Hochberg adjustment is used
only as an exploratory multiple-screening quantity \citep{benjamini1995};
the resulting tiers are observational audit summaries, not labels for
instrument failure, replacement, bias, or pollution correction.
